% archon:physics
% archon:covers ArchonPhysicsConsumers/Thermalization/problem_kinetic_time_rescaling.lean
% archon:source-report reports/thermalization/problem_kinetic_time_rescaling.source.json

\chapter{Generated conditional kinetic rescaling}
\label{ch:ArchonPhysics_Generated_problem_kinetic_time_rescaling}

The theorem concerns a supplied effective ODE and a supplied uniqueness hypothesis. It neither derives a kinetic equation from the microscopic lattice nor proves a microscopic thermalization law. The named conjecture below is only its current, input-dependent formula and is not an asymptotic theorem.

\begin{definition}[Effective kinetic model]
\label{def:kinetic:model}
\lean{ArchonPhysics.Generated.KineticTimeRescaling.EffectiveKineticModel}
A model contains a coupling-independent collision map $C:X\to X$.
\end{definition}
\begin{proof}This record packages precisely the raw collision operator.\end{proof}

\begin{definition}[Kinetic solution]
\label{def:kinetic:solution}
\lean{ArchonPhysics.Generated.KineticTimeRescaling.SolvesKineticEquation}
\uses{def:kinetic:model}
$D$ solves the effective equation when $D(0)=D_0$ and $D'(t)=g^2C(D(t))$.
\end{definition}
\begin{proof}Both initial-value and derivative conditions are included in the predicate.\end{proof}

\begin{definition}[State-event first hit]
\label{def:kinetic:first-hit}
\lean{ArchonPhysics.Generated.KineticTimeRescaling.firstHittingTime}
For an event $P$ and trajectory $D$, take the ENNReal infimum of positive $t$ with $P(D(t))$.
\end{definition}
\begin{proof}The real-valued trajectory is evaluated at the real part of the ENNReal time.\end{proof}

\begin{definition}[Current microscopic formula placeholder]
\label{def:kinetic:microscopic-conjecture}
\lean{ArchonPhysics.Generated.KineticTimeRescaling.microscopicEquilibrationTimeScalingConjecture}
For every $(\lambda,\varepsilon,n)$ satisfying the displayed guards, the current predicate permits a positive coefficient $C$ for an exact inverse-square expression. Its quantifier placement permits $C$ to depend on all inputs; it is therefore not the desired fixed-prefactor asymptotic law.
\end{definition}
\begin{proof}This is an unproved named proposition only; no use of it occurs in the theorem.\end{proof}

\begin{theorem}[Conditional effective solution and hitting-time rescaling]
\label{thm:physics:kinetic_time_rescaling:target}
\lean{ArchonPhysics.Generated.KineticTimeRescaling.kinetic_time_rescaling_physics_formalization_target}
\uses{def:kinetic:model, def:kinetic:solution, def:kinetic:first-hit}
For $g\ne0$, two solutions of the same initial-value problem and an explicit uniqueness principle satisfy $D_g(t)=D_1(g^2t)$; their compatible state-event first hits differ by the inverse ENNReal factor associated to $g^2$.
\end{theorem}
\begin{proof}
Reparameterize the unit-coupling solution by $g^2t$ and apply the chain rule; it satisfies the $g$ equation and same initial condition. Uniqueness gives the trajectory identity. Map the positive event-time set through multiplication by the positive ENNReal factor $g^2$ and transport its infimum through the resulting order isomorphism.
\end{proof}
